% Mục 6: đánh giá định lượng lời giải thích (cấp độ 3). Chỉ dùng val.
% Nguồn số: experiments/xai-metrics-v2/{REPORT.md,summary.json,cases.json} (main_r3, v2) và
% experiments/xai-metrics/REPORT.md (countonly, v1; chỉ những số mà v2 không làm đổi).
\section{Đánh giá định lượng lời giải thích (cấp độ 3)}\label{sec:xai}

Cấp độ 3 yêu cầu đo fidelity (độ trung thực), validity (độ hợp lệ) và sparsity (độ gọn) của lời giải thích trên ít nhất 10 truy vấn.
Mục này mở rộng phép thử xoá của Mục~\ref{sec:l1-cases} ra toàn bộ 10\,633 câu val, 3 seed mỗi mô hình.
Chỉ dùng val: không train, không chấm test, không chỉnh mô hình theo kết quả đo.
Ảnh được giải thích là ảnh mô hình xếp hạng 1 cho mỗi câu, dù đúng hay sai.
Phép đo có hai bản: v1 (\file{experiments/xai-metrics/}) và v2 (\file{experiments/xai-metrics-v2/}), là v1 sau khi soát lại (Mục~\ref{sec:xai-v1v2}).
\texttt{main\_r3} (cấp 1) đã có đủ số v2.
Mô hình chính \texttt{countonly\_factual\_r3} (gọi tắt là countonly) chưa chạy v2: ở đây chỉ dùng số v1 của nó ở các độ đo mà v2 không làm đổi, các ô khác ghi \todo{chờ v2 countonly}.

\subsection{Lời giải thích E và phép sửa đồ thị}\label{sec:xai-setup}

E được đọc ra từ chính lượt xếp hạng (Mục~\ref{sec:method-expl}), không tốn thêm lượt chấm nào.
Mỗi phần $j$ của ảnh hạng 1 (một node vật thể hoặc một bộ ba có nhãn) nhận term
\begin{equation}
  t_j=\frac{\partial s}{\partial S_{\kappa}}\sum_i w_i\,a_{ij}\,\mathit{sim}_{ij},
  \label{eq:xai-term}
\end{equation}
với $\kappa\in\{\text{obj},\text{tri}\}$ là kênh của phần $j$, $S_\kappa$ là điểm kênh và $s$ là điểm cuối.
E là danh sách các phần xếp theo $t_j$ giảm dần; top-$k$ là $k$ phần đầu.
Với fusion raw của countonly, đạo hàm trong \eqref{eq:xai-term} là trọng số kênh ($b$ hoặc $c$ ở Mục~\ref{sec:method-main}), nên các term cộng lại đúng bằng phần đồ thị của điểm cuối.
Với fusion z-score của \texttt{main\_r3}, đạo hàm được tính chính xác bằng autograd trên 50 ứng viên, để node và bộ ba được so trên cùng thang điểm cuối.

Mọi độ đo đều sửa đồ thị của ảnh hạng 1 rồi chấm lại bằng chính mô hình đó (\file{src/explain/xai_metrics.py}).
\emph{Bỏ một node} là bỏ node cùng mọi cạnh chạm vào nó, nên các bộ ba chạm vào nó cũng mất; cạnh kNN không được dựng lại.
\emph{Bỏ một bộ ba} là bỏ đúng cạnh có nhãn đó.
\emph{Chỉ giữ top-$k$} là giữ các phần đó cùng node tổng (nếu có); bộ ba được giữ thì hai node đầu--cuối của nó ở lại làm ``carrier'', vẫn truyền tin qua GAT nhưng không được chấm ở kênh vật thể.
Ảnh sau khi sửa được mã hoá lại bằng cả mô hình, kể cả GAT; 49 ứng viên còn lại giữ nguyên điểm kênh, còn fusion được tính lại trên cả 50 ứng viên.

Mốc \emph{gốc} là ảnh không sửa, chấm lại qua đúng đường chấm đó; ở cả 3 seed của \texttt{main\_r3}, ảnh này vẫn hạng 1 ở mọi câu, sai số điểm tối đa $2{,}0\cdot10^{-5}$.
Mốc \emph{rỗng} bỏ mọi phần; \emph{hiệu ứng đồ thị} là $\Delta=s_{\text{gốc}}-s_{\text{rỗng}}$.
Mốc \emph{ngẫu nhiên} (NN) là trung bình 10 lần bốc cùng số node và cùng số bộ ba với top-$k$, rng cố định theo (seed, ảnh, câu).
Câu có E rỗng (39--45 câu mỗi seed với \texttt{main\_r3}) bị loại khỏi trung bình fidelity.
R@1 tính lại trong lúc đo trùng \file{metrics.json} của từng checkpoint (script kiểm bằng \texttt{assert}).

\subsection{Độ đo}\label{sec:xai-metrics}

\textbf{Fidelity} (tên theo~\cite{xai-yuan2023survey}), với $k\in\{1,3,5\}$.
Fid$^{+}$@$k$ bỏ top-$k$ của E và đếm tỉ lệ câu ảnh mất hạng 1; Fid$^{-}$@$k$ chỉ giữ top-$k$ và ghi hạng của ảnh.
Đường cong deletion / insertion~\cite{xai-petsiuk2018rise} bỏ dần, hoặc chỉ giữ dần, $k=1,\dots,10$ phần đầu; $\mathrm{AUC}=\frac{1}{10}\sum_{k}(s_k-s_{\text{rỗng}})/\Delta$, deletion thấp là tốt, insertion cao là tốt.
AUC có thể âm hoặc lớn hơn 1, vì sửa đồ thị có thể đẩy điểm ra ngoài khoảng $[s_{\text{rỗng}},s_{\text{gốc}}]$.
Fid$_{\alpha}^{\pm}$~\cite{xai-zheng2024fidalpha} bỏ mỗi phần trong 5 phần đầu (Fid$_{\alpha}^{+}$), hoặc mỗi phần còn lại (Fid$_{\alpha}^{-}$), với xác suất $\alpha=0{,}1$, rồi lấy phần $\Delta$ bị mất, trung bình 10 lần bốc.
Đồ thị sau khi sửa vẫn gần đồ thị gốc, và lời giải thích trung thực cho Fid$_{\alpha}^{+}>{}$Fid$_{\alpha}^{-}$.
\emph{Tính cộng} (additivity): vì E là một phép phân rã điểm, tổng term của top-$k$ là ``dự đoán'' của điểm mất khi bỏ top-$k$; nhóm so nó với điểm mất đo được qua tỉ số tổng (tổng đo được chia tổng dự đoán, trên mọi câu) và tương quan Spearman giữa các câu.

\textbf{Validity.}
Về phản thực tế (counterfactual), $k_{\text{flip}}$ là $k\le 10$ nhỏ nhất mà bỏ top-$k$ làm ảnh mất hạng 1.
Về ngữ nghĩa, với mỗi node vật thể của câu, nhóm xét node ảnh có $a_{ij}$ lớn nhất có cùng loại với nó không.
\emph{Strict}: cùng nhãn sau alias VG, hoặc synset của node ảnh trùng hay là hyponym của nghĩa danh từ đầu tiên của từ trong câu; node người chỉ khớp với từ chung (person, people\ldots).
\emph{Lenient} là tiêu chí của v1: cùng nhãn sau alias, hoặc cùng loại theo \texttt{same\_type} (chấp nhận ``woman'' $\sim$ ``man'' và hypernym hai chiều), hoặc cosine CLIP $\ge 0{,}955$.
Dòng bộ ba khớp đủ khi chủ thể và đối tượng khớp strict và vị từ khớp sau alias.
Khớp nhóm số lượng chỉ chấm mô hình có bảng $g$.

\textbf{Sparsity}, theo phần của ảnh: $k_{80}$ là $k\le 10$ nhỏ nhất mà chỉ giữ top-$k$ còn ít nhất 80\% $\Delta$; $n_{80}$ là số phần đầu của E chiếm 80\% khối term dương.

\begin{table*}[t]
  \caption{Kết quả chính trên val, trung bình $\pm$ std qua 3 seed}
  \label{tab:xai-main}
  \centering
  \footnotesize
  \begin{tabular*}{\textwidth}{@{\extracolsep{\fill}}lcccc@{}}
    \toprule
    & \multicolumn{2}{c}{\texttt{main\_r3} (v2)} & \multicolumn{2}{c}{countonly$^{a}$} \\
    \cmidrule{2-3}\cmidrule{4-5}
    Độ đo & E & NN & E & NN \\
    \midrule
    \multicolumn{5}{@{}l}{\emph{Fidelity: bỏ top-$k$ (Fid$^{+}$), \% câu ảnh mất hạng 1 (cao là tốt)}} \\
    \quad $k=1$ & \pmstd{20,5}{1,8} & \pmstd{3,0}{0,2} & \pmstd{34,3}{4,8} & \pmstd{4,4}{0,5} \\
    \quad $k=3$ & \pmstd{34,9}{2,8} & \pmstd{6,5}{0,3} & \pmstd{58,0}{4,1} & \pmstd{9,4}{0,9} \\
    \quad $k=5$ & \pmstd{43,6}{3,1} & \pmstd{9,8}{0,5} & \pmstd{68,0}{3,0} & \pmstd{12,9}{0,7} \\
    \quad Tính cộng: điểm mất / tổng term của top-3 & \pmstd{0,063}{0,016} & & \todo{chờ v2 countonly} & \\
    \multicolumn{5}{@{}l}{\emph{Fidelity: chỉ giữ top-$k$ (Fid$^{-}$), hạng của ảnh (thấp là tốt)}} \\
    \quad $k=1$ & \pmstd{3,00}{0,03} & \pmstd{9,19}{0,55} & \pmstd{4,18}{0,83} & \pmstd{10,13}{2,26} \\
    \quad $k=3$ & \pmstd{2,16}{0,08} & \pmstd{7,38}{0,36} & \pmstd{2,22}{0,74} & \pmstd{8,13}{1,40} \\
    \quad $k=5$ & \pmstd{2,17}{0,06} & \pmstd{6,89}{0,33} & \pmstd{1,86}{0,54} & \pmstd{6,93}{0,96} \\
    \multicolumn{5}{@{}l}{\emph{Validity}} \\
    \quad Câu có $k_{\text{flip}}\le 10$ (\%); trung vị $k_{\text{flip}}$ & \pmstd{60,3}{2,9}; 3 & & \todo{chờ v2 countonly} & \\
    \quad Vật thể khớp strict (\%) & \pmstd{58,3}{0,3} & & \todo{chờ v2 countonly} & \\
    \quad Vật thể khớp lenient (\%) & \pmstd{63,8}{0,2} & & \pmstd{70,0}{1,0} & \\
    \quad Bộ ba khớp đủ chủ thể, đối tượng, vị từ (\%) & \pmstd{5,4}{0,2} & & \todo{chờ v2 countonly} & \\
    \quad Khớp nhóm số lượng (\%) & không chấm$^{b}$ & & \pmstd{66,7}{0,8} & \\
    \multicolumn{5}{@{}l}{\emph{Sparsity (đơn vị: phần của ảnh)}} \\
    \quad $k_{80}$; \% câu đạt trong $k\le 10$ & \pmstd{2,65}{0,14}; 55,7 & & \todo{chờ v2 countonly} & \\
    \quad $n_{80}$; \% số phần của ảnh & \pmstd{15,4}{0,2}; 33,0 & & \todo{chờ v2 countonly} & \\
    \bottomrule
  \end{tabular*}
  \par\smallskip
  {\scriptsize\raggedright
    NN: trung bình 10 lần chọn ngẫu nhiên cùng số node và bộ ba.
    $^{a}$Số v1, chỉ ở các độ đo mà v2 không làm đổi: với fusion raw, E của v1 và v2 trùng nhau; v1 tính thêm 31--36 câu E rỗng mỗi seed (0,3\%) như ``không đổi'' và lấy lượt chấm thường làm gốc (lệch tối đa $1{,}8\cdot10^{-7}$ so với chấm lại).
    Lenient là tiêu chí ngữ nghĩa của v1.
    $^{b}$Chỉ chấm mô hình có bảng số lượng.\par}
  \tabsrc{\file{experiments/xai-metrics-v2/summary.json}, \file{experiments/xai-metrics-v2/REPORT.md} (\texttt{main\_r3}); \file{experiments/xai-metrics/REPORT.md} (countonly)}
\end{table*}

\subsection{Kết quả trên toàn bộ val}\label{sec:xai-results}

Theo Bảng~\ref{tab:xai-main}, ở cả hai mô hình, E vượt mốc ngẫu nhiên theo cả hai chiều của fidelity: bỏ top-3 làm ảnh mất hạng 1 ở số câu gấp 5--6 lần bỏ ngẫu nhiên; chỉ giữ top-3 cho hạng trung bình khoảng 2,2, so với 7--8 của mốc ngẫu nhiên.

Với ảnh hạng 1 sai, $k_{\text{flip}}\le 10$ ở 74,4\% câu; với ảnh hạng 1 đúng, tỉ lệ này là 44,1\%.
Với 5\,325 câu mỗi seed có ảnh hạng 1 sai, bỏ top-5 của E làm ảnh đúng tăng hạng ở 30,7\% câu và lên hạng 1 ở 13,1\% câu.
Về ngữ nghĩa, node ảnh có attention lớn nhất cùng loại với node câu ở 58,3\% dòng vật thể theo tiêu chí strict; với dòng bộ ba, tỉ lệ khớp đủ cả ba thành phần là 5,4\%.

Về sparsity theo hiệu ứng, ở 55,7\% câu, trung bình 2,65 phần đầu giữ được 80\% $\Delta$; ở các câu còn lại, 10 phần đầu chưa đạt mức này.
Theo khối term, cần 15,4 phần (33,0\% số phần của ảnh) để đạt 80\% term dương.

\subsection{Thứ tự và độ lớn của term trong E}\label{sec:xai-findings}

\textbf{Bộ ba đứng đầu E thường không quyết định thứ hạng.}
Với \texttt{main\_r3}, bộ ba đứng đầu E ở 70,5\% câu.
Bỏ top-1 làm ảnh mất hạng 1 ở \pmstd{49,3}{0,3}\% câu khi nó là node (ngẫu nhiên: \pmstd{4,4}{0,1}\%), nhưng chỉ ở \pmstd{8,4}{0,4}\% câu khi nó là bộ ba (ngẫu nhiên: \pmstd{2,5}{0,2}\%).
Countonly (v1) cho cùng xu hướng: bộ ba đứng đầu E ở 44\% câu, và hai con số là \pmstd{46,5}{0,8}\% và \pmstd{19,7}{3,2}\%.
Term của kênh bộ ba phân tán hơn kênh vật thể: ở \texttt{main\_r3}, top-3 bộ ba chiếm \pmstd{22,8}{0,2}\% term dương của kênh, còn top-3 node chiếm \pmstd{63,1}{0,9}\%.

\textbf{Độ lớn của term không dự đoán điểm mất.}
Với \texttt{main\_r3}, tổng term của top-3 là \pmstd{6,12}{0,96}, nhưng bỏ top-3 chỉ làm điểm giảm \pmstd{0,375}{0,036}: tỉ số tổng là \pmstd{0,063}{0,016}, Spearman là \pmstd{$-0{,}035$}{0,016}.
Trên các câu của Mục~\ref{sec:xai-compare}, occlusion (tỉ số \pmstd{1,153}{0,017}, Spearman \pmstd{0,888}{0,011}) và IG (\pmstd{1,434}{0,037}; \pmstd{0,815}{0,009}) cộng tính tốt hơn nhiều.
Với countonly (v1), chỉ giữ top-5 cho 116,6\% $\Delta$, tức điểm cao hơn ảnh gốc.
Vì độ lớn term của E không ước lượng được điểm mất, demo kèm mỗi dòng của E một phép chấm lại khi bỏ phần đó (cột ``bỏ phần này'', Mục~\ref{sec:demo}).

\textbf{Attention của GAT không chỉ ra cạnh quyết định thứ hạng.}
Đề yêu cầu bỏ hoặc nhiễu cả các cạnh dùng làm lời giải thích; ngoài các bộ ba (cạnh có nhãn) trong E, nhóm thử cả các cạnh mà GAT chú ý nhất: với 3 node đầu của E, cắt đường truyền tin của các cạnh vào có attention lớn nhất (giữ bộ ba), rồi so với cắt cùng số cạnh vào ngẫu nhiên của từng node.
Trên 2\,803 câu mỗi seed có nhiều cạnh vào hơn số cạnh cắt, tỉ lệ mất hạng 1 là \pmstd{1,0}{0,6}\% so với \pmstd{0,7}{0,4}\%, và điểm đổi \pmstd{$-0{,}0011$}{0,0007} so với \pmstd{$-0{,}0010$}{0,0005} (\texttt{main\_r3}).
Cắt cạnh có attention cao không làm thay đổi kết quả nhiều hơn cắt cạnh ngẫu nhiên; trọng số attention của GAT vì vậy không xác định được cạnh quyết định thứ hạng, tương tự nhận xét của~\cite{xai-jain2019attention}.
Số tương ứng của countonly: \todo{chờ v2 countonly}.

\subsection{Bản đo v1 sai ở đâu}\label{sec:xai-v1v2}

v1 được soát lại trước khi đưa vào báo cáo; nó giữ nguyên, kèm mục đính chính ở đầu \file{experiments/xai-metrics/REPORT.md}.
v2 sửa sáu điểm.
(1) Với \texttt{main\_r3}, v1 so term thô của hai kênh, dù fusion z-score cân hai kênh khác nhau; v2 thêm đạo hàm ở~\eqref{eq:xai-term}.
Sau khi đưa các term về cùng thang điểm, bộ ba đứng đầu E ở 70,5\% câu thay vì 1\%, và Fid$^{+}$@3 giảm từ \pmstd{65,7}{0,3}\% xuống \pmstd{34,9}{2,8}\%.
(2) v1 lấy lượt chấm thường làm gốc và tính câu E rỗng là ``không đổi''.
(3) Sparsity của v1 chia số dòng của bảng demo cho số phần của ảnh, hai đơn vị khác nhau.
(4) Tiêu chí ngữ nghĩa của v1 chấp nhận cả cặp khác loại; v2 giữ nó dưới tên lenient và thêm strict.
(5) v1 so khớp nhóm số lượng với \texttt{main\_r3} chấm trên đồ thị cấp 1 chưa gộp trùng, nơi hộp trùng làm tăng số đếm.
(6) v1 bỏ cạnh có nhãn là bỏ luôn bộ ba; đo lại dạng này cho 5,7\% câu mất hạng 1 (ngẫu nhiên: 5,1\%), so với 1,0\% khi chỉ cắt truyền tin, nên hiệu ứng mà v1 thấy phần lớn đến từ việc mất bộ ba.

\subsection{So sánh với tám cách giải thích khác}\label{sec:xai-compare}

\begin{table}[t]
  \caption{So sánh 9 cách giải thích, \texttt{main\_r3} (v2)}
  \label{tab:xai-explainers}
  \centering
  \footnotesize
  \setlength{\tabcolsep}{3pt}
  \begin{tabular*}{\columnwidth}{@{\extracolsep{\fill}}lrrrrrr@{}}
    \toprule
    Cách & Del.\,$\downarrow$ & Ins.\,$\uparrow$ & Fid$_{\alpha}^{+}$ & Fid$_{\alpha}^{-}$ & $\rho$ & T5 \\
    \midrule
    Shapley$^{\dagger}$ & $-0{,}683$ & 1,532 & 0,118 & $-0{,}046$ & chuẩn & chuẩn \\
    Occlusion & $-0{,}244$ & 1,234 & 0,106 & $-0{,}046$ & 0,476 & 58,5 \\
    IG & $-0{,}445$ & 0,646 & 0,090 & $-0{,}024$ & 0,288 & 62,9 \\
    GNNExplainer$^{\dagger}$ & $-0{,}169$ & $-0{,}336$ & 0,091 & $-0{,}000$ & 0,255 & 56,7 \\
    \textbf{E} & 0,237 & 0,764 & 0,050 & 0,016 & 0,460 & 31,8 \\
    CLIP & 0,283 & 0,234 & 0,047 & 0,021 & 0,352 & 35,4 \\
    Gradient$\times$input & 0,826 & $-1{,}038$ & 0,012 & 0,057 & $-0{,}133$ & 17,8 \\
    Ngẫu nhiên & 0,838 & $-0{,}621$ & 0,013 & 0,053 & $-0{,}012$ & 14,3 \\
    E đảo ngược & 1,009 & $-2{,}300$ & 0,002 & 0,068 & $-0{,}460$ & 2,1 \\
    \bottomrule
  \end{tabular*}
  \par\smallskip
  {\scriptsize\raggedright
    Val, trung bình 3 seed; std trong file nguồn, lớn nhất ở GNNExplainer (Ins.\ $\pm$\,0,951) và Shapley (Del.\ $\pm$\,0,226).
    Del., Ins.: AUC của đường cong deletion và insertion.
    Fid$_{\alpha}^{\pm}$: 5 phần đầu, $\alpha=0{,}1$.
    $\rho$: Spearman với Shapley; T5: \% top-5 trùng Shapley.
    $^{\dagger}$195 câu mỗi seed; các dòng khác 977 câu.\par}
  \tabsrc{\file{experiments/xai-metrics-v2/REPORT.md}}
\end{table}

Nhóm đo thêm 8 cách giải thích trên cùng các phần của ảnh, bằng cùng các độ đo, trên một mẫu cố định 1\,000 câu val (Bảng~\ref{tab:xai-explainers}): ngẫu nhiên, E đảo ngược, cosine CLIP (không dùng mô hình), gradient $\times$ input, Integrated Gradients~\cite{xai-sundararajan2017ig} (IG, 20 bước), occlusion~\cite{xai-zeiler2014occlusion} (bỏ từng phần, trung bình 46 lượt chấm mỗi câu), GNNExplainer~\cite{xai-ying2019gnnexplainer} và giá trị Shapley lấy mẫu 16 hoán vị~\cite{xai-castro2009shapley}.
Hàm giá trị của Shapley là biên giữa ảnh được giải thích và ứng viên tốt nhất còn lại khi chỉ giữ một tập phần, theo ý tưởng của RankingSHAP~\cite{xai-heuss2025rankingshap}; nó tốn khoảng 740 lượt chấm mỗi câu và được dùng làm chuẩn.
IG thoả completeness với sai số trung bình \pmstd{1,1}{0,1}\% $\Delta$; Shapley thoả efficiency với sai số tối đa $1{,}3\cdot10^{-7}$.
Bốn cách có deletion AUC âm: Shapley, occlusion, IG và GNNExplainer. E có Fid$_{\alpha}^{+}>{}$Fid$_{\alpha}^{-}$ nhưng deletion AUC dương (0,237).
E tốt hơn ngẫu nhiên và gradient $\times$ input ở mọi cột. Deletion AUC của E không khác CLIP có ý nghĩa thống kê (Wilcoxon ghép cặp, $p$ lớn nhất 0,96); mọi cách còn lại khác E với $p<10^{-11}$ ở cả 3 seed.
Tương quan thứ tự giữa E và Shapley (Spearman \pmstd{0,460}{0,007}) gần bằng của occlusion (\pmstd{0,476}{0,007}), trong khi E không cần lượt chấm thêm.
Top-5 của E trùng 31,8\% với Shapley (occlusion: 58,5\%).
E đảo ngược cho kết quả thấp nhất ở mọi cột, nên thứ tự của E chứa thông tin về thứ hạng.

\subsection{Bằng chứng trên 13 truy vấn}\label{sec:xai-cases}

\begin{table}[t]
  \caption{Bằng chứng trên 13 câu val, \texttt{main\_r3} seed 0 (v2)}
  \label{tab:xai-cases}
  \centering
  \footnotesize
  \begin{tabular*}{\columnwidth}{@{\extracolsep{\fill}}llcrrc@{}}
    \toprule
    Câu & Ảnh & Top-3 của E & Bỏ top-3 & Giữ top-3 & $k_{\text{flip}}$ \\
    \midrule
    302038\_1 & Đ & 3 bộ ba & 1 (1,0) & 1 (3,2) & 7 \\
    422698\_1 & Đ & 3 bộ ba & 1 (1,0) & 11 (12,2) & $>10$ \\
    357943\_3 & Đ & 3 bộ ba & 1 (1,0) & 2 (11,7) & 5 \\
    445695\_1 & Đ & 3 node & 11 (2,4) & 1 (8,0) & 1 \\
    574350\_4 & Đ & 3 bộ ba & 1 (1,0) & 2 (2,0) & 7 \\
    1153\_3 & S & 3 bộ ba & 1 (1,0) & 2 (2,8) & 9 \\
    279933\_1 & S & 3 bộ ba & 1 (1,0) & 4 (10,7) & 8 \\
    159627\_2 & Đ & 3 node & 3 (1,0) & 3 (5,8) & 3 \\
    34471\_0 & S & 3 bộ ba & 1 (1,0) & 1 (2,2) & $>10$ \\
    526446\_0 & S & 3 node & 45 (21,9) & 1 (22,0) & 1 \\
    472772\_2 & S & 3 bộ ba & 1 (1,0) & 1 (2,1) & 4 \\
    237041\_4 & S & 3 bộ ba & 1 (1,0) & 1 (5,5) & $>10$ \\
    151938\_1 & S & 3 bộ ba & 1 (1,0) & 1 (2,5) & 7 \\
    \bottomrule
  \end{tabular*}
  \par\smallskip
  {\scriptsize\raggedright
    Ảnh: ảnh hạng 1 đúng (Đ) hay sai (S).
    Hai cột giữa: hạng của ảnh sau khi bỏ, hoặc chỉ giữ, top-3 của E; trong ngoặc là trung bình 10 lần chọn ngẫu nhiên.\par}
  \tabsrc{\file{experiments/xai-metrics-v2/cases.json}}
\end{table}

Bảng~\ref{tab:xai-cases} dùng lại 10 câu của Mục~\ref{sec:l1-cases} và 3 câu ví dụ trong
\file{experiments/improve/explanations/}, lần này chấm bằng \texttt{main\_r3} seed 0.
Ở 3 câu có top-3 của E là node, bỏ top-3 đưa ảnh xuống hạng 11, 3 và 45 (ngẫu nhiên: 2,4; 1,0; 21,9).
Ví dụ, ở câu 526446\_0, top-3 của E là các node catcher, dirt, guard; bỏ ba node này đưa ảnh (sai) xuống hạng 45, còn chỉ giữ chúng thì ảnh vẫn hạng 1.
Ở 10 câu còn lại, top-3 đều là bộ ba, và bỏ cả ba thì ảnh vẫn hạng 1.
Chỉ giữ top-3 thì hạng của ảnh thấp hơn trung bình ngẫu nhiên ở 12/13 câu (câu còn lại bằng nhau), và $k_{\text{flip}}\le 10$ ở 10/13 câu.
Kết quả trên 13 câu cùng chiều với kết quả trên toàn val: khi top-3 của E là node, bỏ chúng làm ảnh mất hạng; khi top-3 là bộ ba, ảnh giữ hạng 1.
